\bibitem{PGCTVI} K. D. Mulmuley, \emph{Geometric Complexity Theory VI: the flip via saturated and positive integer programming}, arXiv:0704.0229v4, supplied revised version. Problems 1.1.1--1.1.4; Definitions 1.2.2, 1.2.4, 1.2.7; Theorems 1.6.1--1.6.3; Hypotheses 1.6.4--1.6.7 and 3.3.2--3.3.4; Lemma 3.1.6; Theorems 3.4.1, 4.1.1, 4.4.1; Sections 4.1.1--4.1.2, equations (4.2)--(4.8). \url{https://arxiv.org/abs/0704.0229v4}.
\bibitem{PGCTVII} K. D. Mulmuley, \emph{Geometric Complexity Theory VII: Nonstandard quantum groups for the plethysm problem}, arXiv:0709.0749. Theorem 1.2; Conjectures 1.3--1.4. \url{https://arxiv.org/abs/0709.0749}.
\bibitem{PGCTVIII} K. D. Mulmuley, \emph{Geometric Complexity Theory VIII: On canonical bases for the nonstandard quantum groups}, arXiv:0709.0751. Conjectures 2.6--2.9 and 3.4--3.7. \url{https://arxiv.org/abs/0709.0751}.
\bibitem{PKM} T. Kahle and M. Micha\l ek, \emph{Obstructions to combinatorial formulas for plethysm}, Electron. J. Combin. 25(1) (2018), P1.41. Theorem 5, Remark 6, Proposition 11, the paragraph after Remark 13, Question 15 and its following paragraph. \url{https://www.combinatorics.org/ojs/index.php/eljc/article/download/v25i1p41/pdf}.
\bibitem{PBOR} E. Briand, R. Orellana and M. Rosas, \emph{Reduced Kronecker coefficients and counter-examples to Mulmuley's strong saturation conjecture SH}, arXiv:0810.3163. Theorem 2.4; Section 2.5; Mulmuley's appendix. \url{https://arxiv.org/abs/0810.3163}.
\bibitem{PIMW} C. Ikenmeyer, K. D. Mulmuley and M. Walter, \emph{On vanishing of Kronecker coefficients}, arXiv:1507.02955. Theorem 1.3. \url{https://arxiv.org/abs/1507.02955}.
\bibitem{PFI} N. Fischer and C. Ikenmeyer, \emph{The computational complexity of plethysm coefficients}, arXiv:2002.00788. Theorems 3.1--3.2. \url{https://arxiv.org/abs/2002.00788}.
\bibitem{PBIP} P. B\"urgisser, C. Ikenmeyer and G. Panova, \emph{No occurrence obstructions in geometric complexity theory}, arXiv:1604.06431. Theorem 1.4. \url{https://arxiv.org/abs/1604.06431}.
\bibitem{PWeilII} P. Deligne, \emph{La conjecture de Weil. II}, Publ. Math. IH\'ES 52 (1980), 137--252. Theorem 3.3.1, Corollaries 3.3.4--3.3.6 and 3.3.9, paragraph 3.3.11. \url{https://www.numdam.org/item/PMIHES_1980__52__137_0.pdf}.
\bibitem{PBBD} A. A. Beilinson, J. Bernstein and P. Deligne, \emph{Faisceaux pervers}, Ast\'erisque 100 (1982). Corollaries 1.4.24--1.4.25, p.~55; equations (4.0.1)--(4.0.2), p.~102; Theorems 5.3.8, 5.4.5, 6.2.5. \url{https://www.sas.rochester.edu/mth/sites/doug-ravenel/otherpapers/BBD-pervers-2x.pdf}.
\bibitem{PHaines} T. J. Haines, \emph{A proof of the Kazhdan--Lusztig purity theorem via the decomposition theorem of BBD}, author note. Theorem 2.0.1, p.~2; Theorem 5.0.4, p.~7. \url{https://math.umd.edu/~tjh/KL_purity1.pdf}.
\bibitem{PMilne} J. S. Milne, \emph{Lectures on \'Etale Cohomology}, version 2.21. Theorem 16.1; Example 16.3; Proposition 18.4(d); Theorem 24.1. \url{https://www.jmilne.org/math/CourseNotes/LEC.pdf}.
\bibitem{PStacks} The Stacks Project, Proposition 59.55.2, tag 03QP, finite-morphism direct images and geometric stalks. \url{https://stacks.math.columbia.edu/tag/03QP}.
\bibitem{PStacksCM} The Stacks Project, Lemma 10.104.2, tag 02JN, the regular-parameter criterion in Cohen--Macaulay local rings. \url{https://stacks.math.columbia.edu/tag/02JN}.
\bibitem{PWeyl} R. Borcherds, M. Haiman, N. Reshetikhin and V. Serganova, notes by T. Johnson-Freyd, \emph{Lectures in Lie Groups and Quantum Groups}, Part I (2016). Example 5.2.0.5 and Proposition 5.2.0.7, pp.~68--69; Theorem 6.1.1.30, pp.~102--103; Theorem 7.6.1.3, p.~167. \url{https://math.berkeley.edu/~serganov/261/part1.pdf}.
\bibitem{PDiscriminant} C. McCrory, T. Shifrin and R. Varley, \emph{Siegel modular forms generated by invariants of cubic hypersurfaces}, arXiv:alg-geom/9403005v3. Section 3.1 and Section 3.1.2(i), printed p.~15: discriminant invariance, degree and singular-hypersurface locus. \url{https://arxiv.org/pdf/alg-geom/9403005}.
\bibitem{PBuse} L. Bus\'e and J.-P. Jouanolou, \emph{A computational approach to the discriminant of homogeneous polynomials}, arXiv:1210.4697v1. Definition 4.6 and Proposition 4.7(i), p.~40; Example 4.12, p.~42; Proposition 4.13, pp.~42--43. \url{https://arxiv.org/pdf/1210.4697}.
\bibitem{PUniversalDenom} H. Derksen, \emph{Universal denominators of Hilbert series}, J. Algebra 285 (2005), 586--607. Definition 1.5; Theorem 1.10; Example 4.7, pp.~597--598. \url{https://sites.lsa.umich.edu/hderksen/wp-content/uploads/sites/614/2018/05/A.I.a.19.pdf}.
\bibitem{PMacdonald} I. G. Macdonald, \emph{Symmetric Functions and Hall Polynomials}, second edition, Oxford, 1995. I (7.3), (7.6)(ii), (7.9); I \S5 Example 2; Appendix A (7.3)--(7.4). \url{https://math.berkeley.edu/~corteel/MATH249/macdonald.pdf}.
\bibitem{PBorger} J. Borger, \emph{$\Lambda$-rings and the field with one element}, arXiv:0906.3146, supplied file \texttt{lrfoe13.tex}. Section 1.1, paragraph on the right adjoint of the big Witt functor. \url{https://arxiv.org/abs/0906.3146}.
\bibitem{PSGA6} P. Berthelot, A. Grothendieck and L. Illusie, eds., \emph{Th\'eorie des intersections et th\'eor\`eme de Riemann--Roch}, SGA 6, LNM 225, 1971. Expos\'e V, Definitions 2.1, 2.4; Expos\'e VI, Theorem 3.2. Original-page retrieval was not completed in this run; the check register distinguishes this background citation from directly inspected primary pages.
\bibitem{PBouc} S. Bouc and K. R\"okaeus, \emph{A note on the $\lambda$-structure on the Burnside ring}. Section 1 and Theorem 1.1, equation (1.2). \url{https://www.lamfa.u-picardie.fr/bouc/Burnside2.pdf}.
\bibitem{PBaez} J. Baez, J. Moeller and T. Trimble, \emph{Schur Functors and Categorified Plethysm}, arXiv:2106.00190. Theorem 6.11. \url{https://arxiv.org/abs/2106.00190}.
\bibitem{PZib1} M. Zibrowius, \emph{Symmetric representation rings are $\lambda$-rings}, New York J. Math. 21 (2015), 1055--1092. Main theorem, p.~1056; Theorem 3.5. \url{https://nyjm.albany.edu/j/2015/21-49v.pdf}.
\bibitem{PZib2} M. Zibrowius, \emph{The gamma-filtration on the Witt ring of a scheme}, arXiv:1411.3198v2. Theorem 2; Introduction, p.~2. \url{https://arxiv.org/abs/1411.3198}.
\bibitem{PPPS} I. Pak, G. Panova and J. P. Swanson, \emph{A combinatorial interpretation for certain plethysm and Kronecker coefficients}, arXiv:2511.02312v3, July 2026 version. Theorems 1.1--1.2 and Lemma 2.1. \url{https://arxiv.org/abs/2511.02312}.
\bibitem{PArchive} A. Grothendieck, supplied manuscript transcriptions, cotes/scans 133/6, 133/8, 133/9, 133/19, 133/28, 2/17, 43/37, 119/217, 12/41, 16/38, 4/7, 79/51 and 89/18; filenames and hashes in the source manifest. The associated supplied study \texttt{SA17\_STUDY\_cote133\_lambda.md} is treated separately from the transcriptions and is corrected where PL-01 applies.
